% !TEX program = xelatex
% Lernplatt - by Can Siebert
% Kurs 22 (Bo 143) - Tag 7 - Lehrskript
% Plattform-Verkauf meistern - SEO, Retention & Bewertungsmanagement
%
% Self-contained xelatex source. Kein biblatex, keine externen Bilder.

\documentclass[11pt,a4paper,oneside]{article}

% --- Encoding & Sprache --------------------------------------------------
\usepackage{polyglossia}
\setdefaultlanguage{german}

% --- Geometrie -----------------------------------------------------------
\usepackage[a4paper,top=2cm,bottom=2cm,outer=2.5cm,inner=2.5cm,bindingoffset=1.5cm]{geometry}

% --- Fonts ---------------------------------------------------------------
\usepackage{fontspec}
\defaultfontfeatures{Ligatures=TeX}

\IfFontExistsTF{Charter}{%
  \setmainfont{Charter}[%
    Numbers=OldStyle,%
    UprightFeatures   = {SmallCapsFeatures={Letters=SmallCaps}}%
  ]%
}{%
  \IfFontExistsTF{Bitstream Charter}{%
    \setmainfont{Bitstream Charter}%
  }{%
    \setmainfont{TeX Gyre Schola}%
  }%
}

\IfFontExistsTF{Source Sans 3}{%
  \setsansfont{Source Sans 3}%
}{%
  \IfFontExistsTF{Source Sans Pro}{%
    \setsansfont{Source Sans Pro}%
  }{%
    \setsansfont{TeX Gyre Heros}%
  }%
}

\newcommand{\caveatfontfallback}{\itshape}
\IfFontExistsTF{Caveat}{%
  \newfontfamily\caveatfont{Caveat}[Scale=1.15]%
}{%
  \newcommand\caveatfont{\caveatfontfallback}%
}

\setmonofont{Latin Modern Mono}

% --- Mikro-Typografie ----------------------------------------------------
\usepackage{microtype}
\linespread{1.15}

% --- Farben & Boxen ------------------------------------------------------
\usepackage{xcolor}
\definecolor{lpInk}{RGB}{20,28,38}
\definecolor{lpAccent}{RGB}{22,90,140}
\definecolor{lpAccentLight}{RGB}{210,228,240}
\definecolor{lpAmber}{RGB}{222,138,30}
\definecolor{lpAmberLight}{RGB}{255,243,217}
\definecolor{lpAmberBorder}{RGB}{196,118,18}
\definecolor{lpGray}{RGB}{96,104,114}
\definecolor{lpGrayLight}{RGB}{238,240,243}

\usepackage[most]{tcolorbox}
\tcbuselibrary{breakable,skins}

\newtcolorbox{eselsbox}[1][Eselsbrücke]{%
  enhanced, breakable,
  colback=lpAmberLight,
  colframe=lpAmberBorder,
  boxrule=0.6pt,
  arc=2pt,
  borderline={0.4pt}{0pt}{lpAmberBorder, dashed},
  left=10pt,right=10pt,top=8pt,bottom=8pt,
  fonttitle=\sffamily\bfseries\color{lpAmberBorder},
  coltitle=lpAmberBorder,
  title={#1},
  attach boxed title to top left={xshift=8pt,yshift=-8pt},
  boxed title style={size=small, colback=lpAmberLight, colframe=lpAmberBorder, boxrule=0.4pt, arc=1pt},
  top=14pt
}

\newtcolorbox{merkbox}[1][Merksatz]{%
  enhanced, breakable,
  colback=lpAccentLight,
  colframe=lpAccent,
  boxrule=0.6pt,
  arc=2pt,
  left=10pt,right=10pt,top=8pt,bottom=8pt,
  fonttitle=\sffamily\bfseries\color{white},
  coltitle=white,
  title={#1},
  attach boxed title to top left={xshift=8pt,yshift=-8pt},
  boxed title style={size=small, colback=lpAccent, colframe=lpAccent, boxrule=0.4pt, arc=1pt},
  top=14pt
}

% --- Listen --------------------------------------------------------------
\usepackage{enumitem}
\setlist{itemsep=2pt, topsep=4pt, parsep=0pt}
\setlist[itemize,1]{label={\textbullet},leftmargin=1.6em}
\setlist[itemize,2]{label={\textendash},leftmargin=1.4em}
\setlist[enumerate,1]{label=\arabic*., leftmargin=1.8em}

% --- Tabellen ------------------------------------------------------------
\usepackage{tabularx}
\usepackage{array}
\usepackage{booktabs}
\usepackage{longtable}

% --- Kopf-/Fusszeilen ----------------------------------------------------
\usepackage{fancyhdr}
\usepackage{lastpage}
\pagestyle{fancy}
\fancyhf{}
\renewcommand{\headrulewidth}{0.3pt}
\renewcommand{\footrulewidth}{0pt}
\fancyhead[L]{\footnotesize\sffamily\color{lpGray}Lernplatt \textperiodcentered{} Kurs 22 \textperiodcentered{} Tag 7}
\fancyhead[R]{\footnotesize\sffamily\color{lpGray}Plattform-SEO, Content \& Bewertungen}
\fancyfoot[C]{\footnotesize\sffamily\color{lpGray}\thepage{} / \pageref{LastPage}}

\fancypagestyle{plain}{%
  \fancyhf{}%
  \renewcommand{\headrulewidth}{0pt}%
  \fancyfoot[C]{\footnotesize\sffamily\color{lpGray}\thepage{} / \pageref{LastPage}}%
}

% --- Section-Styling -----------------------------------------------------
\usepackage[explicit]{titlesec}

\titleformat{\section}[block]
  {\sffamily\Large\bfseries\color{lpAccent}}
  {\thesection.}{0.6em}{#1}
  [{\vspace{2pt}\color{lpAccent}\titlerule[0.6pt]}]

\titleformat{\subsection}[block]
  {\sffamily\large\bfseries\color{lpInk}}
  {\thesubsection}{0.6em}{#1}

\titleformat{\subsubsection}[block]
  {\sffamily\normalsize\bfseries\color{lpInk}}
  {}{0pt}{#1}

\titlespacing*{\section}{0pt}{14pt}{8pt}
\titlespacing*{\subsection}{0pt}{10pt}{4pt}
\titlespacing*{\subsubsection}{0pt}{8pt}{2pt}

% --- Hyperlinks ----------------------------------------------------------
\usepackage[hidelinks,unicode]{hyperref}

% --- TOC -----------------------------------------------------------------
\renewcommand{\contentsname}{\sffamily\Large\bfseries\color{lpAccent}Inhalt}

% --- Helfer --------------------------------------------------------------
\newcommand{\akzent}[1]{{\caveatfont\color{lpAmber}#1}}
\newcommand{\fachbegriff}[1]{\textbf{#1}}
\newcommand{\quelle}[1]{{\footnotesize\color{lpGray}(#1)}}

% =========================================================================
\begin{document}

% --- COVER ---------------------------------------------------------------
\begin{titlepage}
  \pagestyle{empty}
  \vspace*{1.5cm}
  \noindent{\sffamily\footnotesize\color{lpGray}LERNPLATT \textperiodcentered{} BY CAN SIEBERT}\\[2pt]
  \noindent{\color{lpAccent}\rule{\linewidth}{1.2pt}}\\[4pt]
  \noindent{\sffamily\footnotesize\color{lpGray}MODUL 4 \textperiodcentered{} TAG 7 \textperiodcentered{} KURS 22 (Bo 143) \textperiodcentered{} 8.\,Juni 2026}

  \vspace{3cm}
  {\sffamily\fontsize{30}{34}\selectfont\bfseries\color{lpInk}Plattform-SEO,\\Content\\\& Bewertungen\par}

  \vspace{0.7cm}
  {\sffamily\large\color{lpGray}Sichtbarkeit, Vertrauen und Wiederkauf\\auf digitalen Plattformen steuern.\par}

  \vspace{1.5cm}
  {\caveatfont\LARGE\color{lpAmber}Lehrskript -- Tag 7 von 16}\par

  \vfill

  \noindent\begin{minipage}{0.55\linewidth}
    {\sffamily\footnotesize\color{lpGray}Lernpfad:}\\
    {\sffamily\small\color{lpInk}Wissen \textrightarrow{} Anwendung \textrightarrow{} Management}\\[10pt]
    {\sffamily\footnotesize\color{lpGray}Bearbeitungsdauer:}\\
    {\sffamily\small\color{lpInk}1 Kurstag (8 Unterrichtseinheiten)}
  \end{minipage}\hfill
  \begin{minipage}{0.4\linewidth}
    \raggedleft
    {\sffamily\footnotesize\color{lpGray}Dozent:}\\
    {\sffamily\small\color{lpInk}Can Siebert}\\[10pt]
    {\sffamily\footnotesize\color{lpGray}Format:}\\
    {\sffamily\small\color{lpInk}Theorie \textperiodcentered{} Fallstudie \textperiodcentered{} Coaching}
  \end{minipage}

  \vspace{0.5cm}
  \noindent{\color{lpAccent}\rule{\linewidth}{0.6pt}}
\end{titlepage}

% --- LERNZIELE -----------------------------------------------------------
\section*{Lernziele dieses Tages}
\addcontentsline{toc}{section}{Lernziele dieses Tages}
\thispagestyle{plain}

\noindent Mit Tag~7 wechselt der Fokus auf den \emph{Plattform-Verkauf}. Wer auf einem Marktplatz oder einer großen Plattform verkauft, hat einen anderen Hebel als auf der eigenen Website: nicht persönliche Beratung, sondern Sichtbarkeit, Content, Bewertungen und Wiederkauf. Heute lernen Sie, diese Hebel strategisch zu nutzen -- und ihre Grenzen zu kennen.

\subsection*{L1 -- Wissen (Reproduktion und Einordnung)}
\begin{itemize}
  \item Grundlagen von Plattform-SEO verstehen.
  \item Sichtbarkeitsfaktoren auf Plattformen kennen.
  \item Keyword-Strategien und Content-Optimierung einordnen.
  \item Bedeutung von Produktdaten, Titeln, Beschreibungen, Bildern und Bewertungen verstehen.
  \item Grundlagen von Kundenretention und Bewertungsmanagement kennen.
\end{itemize}

\subsection*{L2 -- Anwendung (Transfer auf konkrete Situationen)}
\begin{itemize}
  \item Plattformangebote hinsichtlich Sichtbarkeit und Verkaufswirkung analysieren.
  \item Keywords, Produkttitel und Beschreibungen zielgerichtet optimieren.
  \item Maßnahmen zur Kundenbindung auf Plattformen entwickeln.
  \item Bewertungsmanagement strukturiert anwenden.
  \item Zielkonflikte zwischen Sichtbarkeit, Markenwirkung, Preispositionierung und Kundenerwartung bewerten.
\end{itemize}

\subsection*{L3 -- Management (Entscheidung und Verantwortung)}
\begin{itemize}
  \item Plattform-SEO und Bewertungsmanagement als strategische Steuerungsinstrumente bewerten.
  \item Entscheidungen unter Unsicherheit, Wettbewerbsdruck und begrenztem Budget treffen.
  \item Verantwortung für Sichtbarkeitsarchitektur, Reputation und Kundenbindung übernehmen.
  \item Plattform-Verkauf nicht nur operativ, sondern als Teil der digitalen Vertriebsstrategie steuern.
\end{itemize}

\begin{merkbox}[Roter Faden]
Sichtbarkeit ist auf Plattformen \emph{ein Produkt}: sie wird über Algorithmen, Bewertungen, Content-Qualität und bezahlte Platzierung verteilt. Wer das nicht versteht, optimiert isolierte Stellschrauben. Wer es versteht, baut eine \textbf{Plattform-Performance-Architektur} -- ein zusammenhängendes System aus Sichtbarkeit, Vertrauen, Conversion und Retention.
\end{merkbox}

\clearpage

% --- 1. EINSTIEG ---------------------------------------------------------
\section{Einstieg: Sichtbarkeit ist nicht Reichweite}

Auf einer Plattform wie Amazon, Otto, eBay oder einem B2B-Marktplatz sind potenziell Millionen Kunden unterwegs. Reichweite \emph{könnte} riesig sein. Tatsächlich aber wird sie über den Algorithmus verteilt -- und der bevorzugt jene Anbieter, die bestimmte Signale liefern: relevante Keywords, vollständige Produktdaten, hohe Klickrate, hohe Conversion, gute Bewertungen, schnelle Lieferung. Wer diese Signale nicht aktiv steuert, bleibt unsichtbar -- selbst mit dem besten Produkt. \quelle{Enge et al., 2015; Fishkin, 2018}

\subsection{Drei Wahrheiten über Plattform-Verkauf}

\begin{enumerate}
  \item \fachbegriff{Sichtbarkeit ist verteilt, nicht garantiert.} Das beste Produkt wird ohne saubere Plattform-SEO schlicht nicht gefunden.
  \item \fachbegriff{Vertrauen ist quantifiziert.} Bewertungssterne, Anzahl Bewertungen, Verkäufer-Bewertung sind nicht Zierde -- sie sind Ranking- und Conversion-Faktoren.
  \item \fachbegriff{Kunden vergleichen brutal.} Auf der Produktseite sehen sie nicht nur uns, sondern drei Wettbewerber. Wer die Conversion nicht aktiv gestaltet, verliert.
\end{enumerate}

\subsection{Was Tag~7 von Tag~5--6 unterscheidet}

Tag 5 hat die kundenzentrierte Kanalwahl behandelt. Tag 6 hat die Omnichannel-Operationalisierung in die eigene Organisation gebracht. Tag 7 wechselt auf die \emph{Außenseite einer fremden Plattform}: Wir verkaufen \emph{innerhalb} eines Systems, dessen Regeln wir nicht setzen, dessen Algorithmus wir nicht kennen, dessen Daten wir nur teilweise sehen. Hier gelten andere Hebel als auf der eigenen Website.

\begin{eselsbox}[Eselsbrücke: \akzent{S-V-K} -- drei Hebel auf der Plattform]
Drei Buchstaben, die jede Plattform-Performance erklären:\\[2pt]
\textbf{S}ichtbarkeit -- werden wir gefunden?\\
\textbf{V}ertrauen -- vertraut der Kunde uns genug zum Kauf?\\
\textbf{K}auf \& Wiederkauf -- konvertieren wir, halten wir den Kunden?\\[2pt]
\emph{``S-V-K'' -- drei Hebel, ein System.}
\end{eselsbox}

% --- 2. PLATTFORM-SEO ---------------------------------------------------
\section{Plattform-SEO: was anders ist}

Klassische Suchmaschinenoptimierung (SEO für Google) und Plattform-SEO (z.\,B.\ Amazon, Marktplätze) teilen viele Prinzipien -- aber nicht alle. Eine Verwechslung kostet Geld. \quelle{vgl.\ Enge et al., 2015}

\subsection{Klassische SEO vs.\ Plattform-SEO}

\begin{tabularx}{\linewidth}{@{}lXX@{}}
\toprule
& \textbf{Klassische SEO} & \textbf{Plattform-SEO} \\
\midrule
\textbf{Suchabsicht} & meist informativ & meist kaufnah \\
\textbf{Ranking-Faktoren} & Inhalt, Backlinks, technische Qualität & Produkttitel, Verkaufszahlen, Bewertungen, Conversion \\
\textbf{Content-Format} & Texte, Blog, Knowledge & Produkttitel, Produktdaten, Bilder, FAQs \\
\textbf{Wirkungshorizont} & Wochen bis Monate & Tage bis Wochen \\
\textbf{Margenwirkung} & indirekt (Leads) & direkt (Verkauf) \\
\bottomrule
\end{tabularx}

\subsection{Die typischen Sichtbarkeitsfaktoren auf Plattformen}

\begin{itemize}
  \item \textbf{Suchbegriff-Relevanz.} Wie genau passen Produkttitel und Beschreibung zum Suchbegriff?
  \item \textbf{Produkttitel.} Klar, suchorientiert, mit Nutzenargument -- nicht nur Modellnummer.
  \item \textbf{Produktbeschreibung.} Struktur, Vollständigkeit, Anwendungsfälle, FAQ.
  \item \textbf{Kategorie.} Richtige Einordnung -- falsche Kategorie killt Sichtbarkeit.
  \item \textbf{Bilder \& Medien.} Anzahl, Qualität, Anwendungssituationen, Videos.
  \item \textbf{Lieferfähigkeit.} Verfügbarkeit, Lieferzeit, Versandkosten.
  \item \textbf{Klickrate.} Wie oft wird unser Listing angeklickt, wenn es angezeigt wird?
  \item \textbf{Conversion Rate.} Wie viele Klicks führen zum Kauf?
  \item \textbf{Retourenquote.} Hohe Rückläufer verschlechtern Sichtbarkeit.
  \item \textbf{Bewertungen.} Anzahl, Durchschnitt, Aktualität.
  \item \textbf{Verkäuferperformance.} Reklamationsquote, Antwortzeiten, Konfliktlösung.
\end{itemize}

\subsection{Risiken reiner Sichtbarkeitsoptimierung}

Wer nur auf Sichtbarkeit optimiert, ohne Conversion und Markenwirkung mitzudenken, kauft sich Probleme ein:

\begin{itemize}
  \item Falsche Kundenerwartungen \textrightarrow{} Rückläufer und schlechte Bewertungen.
  \item Preisdruck \textrightarrow{} dauerhafte Margenerosion.
  \item Margenverlust durch teure Werbeplätze, die nicht refinanziert werden.
  \item Verwässerung der Markenpositionierung.
\end{itemize}

\begin{merkbox}[Sichtbarkeits-Faustregel]
Sichtbarkeit ohne Conversion ist ein teures Hobby. Bevor Sie Budget in Reichweitenwerbung pumpen, prüfen Sie: ist die Produktseite überhaupt geeignet, einen Klick in einen Kauf zu wandeln? Wenn nicht, optimieren Sie die Conversion zuerst.
\end{merkbox}

% --- 3. KEYWORD-STRATEGIE -----------------------------------------------
\section{Keyword-Strategie}

Eine sinnvolle Plattform-SEO-Strategie beginnt bei den Suchbegriffen der Zielgruppe -- nicht bei den eigenen Produktbezeichnungen. ``TX 18V Set'' ist eine Modellnummer; \emph{niemand} sucht nach ihr ohne Vorkenntnisse. Kunden suchen nach Anwendungsproblemen und Nutzenversprechen.

\subsection{Fünf Keyword-Arten}

\begin{enumerate}
  \item \fachbegriff{Hauptkeywords (Generika).} Kurz, hohes Volumen, hoher Wettbewerb -- z.\,B.\ ``Akkuschrauber''.
  \item \fachbegriff{Nebenkeywords.} Etwas spezifischer, mittleres Volumen -- z.\,B.\ ``Akkuschrauber 18V''.
  \item \fachbegriff{Long-Tail-Keywords.} Längere, präzisere Phrasen, niedrigeres Volumen, höhere Kaufabsicht -- z.\,B.\ ``Profi-Akkuschrauber 18V mit Ersatzakku für Handwerker''.
  \item \fachbegriff{Problemorientierte Keywords.} Beschreiben die Anwendungssituation -- z.\,B.\ ``Akkuschrauber für Dauereinsatz Montage''.
  \item \fachbegriff{Kaufnahe Keywords.} Mit Kauf-Signal -- z.\,B.\ ``Akkuschrauber kaufen Werkstatt''.
\end{enumerate}

\subsection{Suchintentionen verstehen}

Jeder Suchbegriff trägt eine \fachbegriff{Suchintention}:

\begin{itemize}
  \item \textbf{Informieren} -- ``Was ist ein Akkuschrauber-Drehmoment?''
  \item \textbf{Vergleichen} -- ``Akkuschrauber Test 2026''.
  \item \textbf{Kaufen} -- ``Profi-Akkuschrauber 18V Set kaufen''.
  \item \textbf{Wiederbestellen} -- ``Ersatzakku Modell TX 18V''.
\end{itemize}

Wer auf einer Plattform Reichweite \emph{ohne} Kaufabsicht einkauft, zahlt für Klicks ohne Umsatz. Auf einem Marktplatz gilt: priorisieren Sie kaufnahe Keywords mit guter Zielgruppenpassung -- nicht die mit dem höchsten Suchvolumen.

\subsection{Keyword-Priorisierung -- nicht nur Suchvolumen}

Eine Keyword-Auswahl wird auf vier Achsen bewertet:

\begin{enumerate}
  \item \textbf{Zielgruppenpassung.} Sucht meine Wunschzielgruppe so?
  \item \textbf{Suchvolumen.} Wie groß ist die Nachfrage?
  \item \textbf{Wettbewerb.} Wie schwer ist die Platzierung?
  \item \textbf{Markenpassung.} Passt der Begriff zur Positionierung -- oder zieht er falsche Kunden?
\end{enumerate}

\subsection{Wann ``billig'' tabu ist}

Ein klassisches Beispiel: das Keyword ``billiger Industriesauger'' hat hohes Volumen, aber niedrige Markenpassung für einen Qualitätsanbieter. Hier ist die Senior-Entscheidung: \emph{bewusst verzichten}, auch wenn Reichweite verlockend wäre. Wer auf ``billig'' optimiert, zieht ``billig''-erwartende Kunden an -- und produziert schlechte Bewertungen.

\begin{eselsbox}[Eselsbrücke: \akzent{Z-S-W-M} -- vier Achsen jeder Keyword-Wahl]
Vier Buchstaben, vier Achsen:\\[2pt]
\textbf{Z}ielgruppenpassung -- sucht der richtige Mensch so?\\
\textbf{S}uchvolumen -- ist die Nachfrage groß genug?\\
\textbf{W}ettbewerb -- können wir hier ranken?\\
\textbf{M}arkenpassung -- zieht das Keyword die richtigen Kunden an?\\[2pt]
\emph{Wer nur ``S'' optimiert, kauft Klicks, keine Umsätze.}
\end{eselsbox}

% --- 4. PRODUKTTITEL & BESCHREIBUNG --------------------------------------
\section{Produkttitel und Produktbeschreibung}

Der \fachbegriff{Produkttitel} ist auf Plattformen die wichtigste einzelne Optimierungsstelle. Er entscheidet über Sichtbarkeit (Ranking), Klickrate (warum klickt der Kunde?) und Erwartung (was kauft er gleich?).

\subsection{Aufbau eines verkaufsstarken Produkttitels}

Ein wirksamer Produkttitel kombiniert:

\begin{enumerate}
  \item \textbf{Produkttyp + Hauptkeyword.} Was ist das eigentlich? (``Profi-Akkuschrauber 18V Set'')
  \item \textbf{Differenzierungsmerkmal.} Wodurch unterscheidet sich das Produkt? (``mit 2 Akkus, Schnellladegerät, robustem Koffer'')
  \item \textbf{Anwendungsbezug.} Für wen / wofür? (``für Handwerk \& Montage'')
\end{enumerate}

\subsubsection{Beispiel}

Vorher: ``Akkuschrauber TX 18V Set''. Nachher: ``Profi-Akkuschrauber 18V Set für Handwerk \& Montage -- mit 2 Akkus, Schnellladegerät und robustem Transportkoffer''. Der zweite Titel rankt besser, wirkt vertrauenswürdiger und reduziert Rückläufer durch unklare Erwartungen.

\subsection{Struktur einer Produktbeschreibung}

\begin{enumerate}
  \item \textbf{Nutzen-Einstieg.} Für welche Anwendung ist das Produkt geeignet? Welches Problem löst es?
  \item \textbf{Wichtigste Vorteile.} 3--5 Punkte (Leistung, Akkulaufzeit, Robustheit, Ergonomie, Zubehör).
  \item \textbf{Technische Daten.} Übersichtlich, vollständig, vergleichbar.
  \item \textbf{Kompatibilität.} Welche Akkus, Bits, Zubehörteile passen? (Reduziert Rückläufer und Beschwerden enorm.)
  \item \textbf{Anwendungsfälle.} Konkrete Einsatzszenarien (Montage, Innenausbau, Werkstatt, Serviceeinsatz).
  \item \textbf{Lieferumfang.} Transparent aufgeführt.
  \item \textbf{FAQ.} Häufige Fragen vor dem Kauf beantworten (Ersatzakku, Ladezeit, Garantie, Versand).
\end{enumerate}

\subsection{Bilder, Videos, Vergleichstabellen}

\begin{itemize}
  \item \textbf{Erstes Bild.} Klar, hoher Kontrast, Produkt im Vordergrund.
  \item \textbf{Anwendungsbilder.} Produkt in realer Einsatzsituation -- erzeugt Identifikation.
  \item \textbf{Vergleichstabellen.} Zeigen Unterschiede zu eigenen Produktvarianten (Basis vs.\ Premium).
  \item \textbf{Kurzvideos.} Sehr wirksam, sofern Plattform sie zulässt.
\end{itemize}

\begin{merkbox}[Content-Faustregel]
Die Produktseite ist die wichtigste Conversion-Fläche im Plattform-Verkauf. Sie ersetzt das Beratungsgespräch. Wenn sie weniger leistet als eine 5-Minuten-Beratung, lassen Sie potentielle Umsätze auf dem Tisch.
\end{merkbox}

% --- 5. KUNDENRETENTION AUF PLATTFORMEN ---------------------------------
\section{Kundenretention auf Plattformen}

Auf Plattformen ist Retention schwerer als auf der eigenen Website, weil die Plattform die Kundenbeziehung ``hält''. Trotzdem gibt es klare Hebel. \quelle{Reichheld, 2003; Lemon \& Verhoef, 2016; Ellis \& Brown, 2017}

\subsection{Was Retention konkret bedeutet}

\begin{itemize}
  \item Der Kunde \textbf{wiederbestellt} bei uns -- statt beim Wettbewerber.
  \item Der Kunde \textbf{erweitert} den Kauf um Zubehör, Bundles, Folgeprodukte.
  \item Der Kunde \textbf{abonniert} regelmäßige Lieferungen oder Verträge.
  \item Der Kunde \textbf{empfiehlt} weiter über Bewertungen oder Mund-zu-Mund.
\end{itemize}

\subsection{Vier Retention-Hebel im Plattform-Verkauf}

\begin{enumerate}
  \item \fachbegriff{Zubehör- und Ersatzteil-Strategie.} Hauptprodukt + Zubehör + Ersatz = wiederkehrender Bedarf. Wer Akkuschrauber verkauft, sollte Akkus, Bits und Ladegeräte sichtbar mitführen -- als Bundle oder als ``Häufig zusammen gekauft''.
  \item \fachbegriff{Bundles.} Geräte- + Zubehör- + Service-Bundles reduzieren Vergleichbarkeit und erhöhen den Warenkorbwert.
  \item \fachbegriff{Nachkaufimpulse.} Sofern die Plattform es zulässt: E-Mail-Strecken nach dem Kauf, Erinnerungen für Verbrauchsmaterial, Empfehlungen für Folgeprodukte.
  \item \fachbegriff{Wiederkaufrate als KPI.} Wer Wiederkauf nicht misst, optimiert ihn nicht. Eine Wiederkaufrate von 25\,\% statt 15\,\% verändert das Geschäftsmodell signifikant.
\end{enumerate}

\subsection{Lieferung, Service, Kommunikation}

Im Plattform-Verkauf entscheidet jenseits des Produkts die operative Qualität:

\begin{itemize}
  \item Lieferzeit eingehalten? Versandkosten transparent?
  \item Anfragen schnell beantwortet?
  \item Reklamationen schmerzfrei abgewickelt?
\end{itemize}

Diese drei Punkte sind keine ``Hygiene'' -- sie sind die Hälfte der Bewertungen, die ein Kunde am Ende abgibt.

% --- 6. BEWERTUNGSMANAGEMENT ---------------------------------------------
\section{Bewertungsmanagement}

Bewertungen sind auf Plattformen \emph{Mehrfach-Asset}: Vertrauenssignal, Ranking-Faktor, Conversion-Treiber und Frühwarnsystem in einem. Wer Bewertungen passiv abwartet, verschenkt einen der wichtigsten Hebel. \quelle{Chevalier \& Mayzlin, 2006; Reichheld, 2003}

\subsection{Warum Bewertungen wirken}

\begin{enumerate}
  \item \textbf{Vertrauen.} Andere Käufer wirken glaubwürdiger als der Anbieter selbst.
  \item \textbf{Ranking.} Plattformalgorithmen bevorzugen besser bewertete Anbieter.
  \item \textbf{Conversion.} Produkte mit > 4{,}5 Sternen werden deutlich häufiger gekauft als solche mit 3{,}8 Sternen -- bei sonst gleichem Angebot.
  \item \textbf{Frühwarnung.} Wiederkehrende Kritik in Bewertungen zeigt operative oder qualitative Probleme an, oft schneller als interne Reports.
\end{enumerate}

\subsection{Vier Aufgaben des Bewertungsmanagements}

\begin{enumerate}
  \item \fachbegriff{Ermöglichen.} Zufriedene Kunden aktiv und regelkonform zur Bewertung anregen -- ohne Anreize, die die Plattform-Regeln verletzen.
  \item \fachbegriff{Analysieren.} Bewertungen regelmäßig (mind.\ wöchentlich) auswerten und nach Themen clustern: Produkt, Lieferung, Service, Erwartung.
  \item \fachbegriff{Beantworten.} Negative Bewertungen sachlich, lösungsorientiert und transparent beantworten -- innerhalb definierter Reaktionszeit.
  \item \fachbegriff{Lernen.} Wiederkehrende Kritikpunkte an Produktmanagement, Logistik oder Content zurückspielen -- mit klarer Owner-Rolle.
\end{enumerate}

\subsection{Umgang mit negativen Bewertungen}

\begin{itemize}
  \item \textbf{Sachlich bleiben.} Persönlich werden lesen 100 künftige Kunden mit.
  \item \textbf{Lösung anbieten.} Konkreter Vorschlag, kein leerer ``Wir nehmen Ihr Feedback ernst''.
  \item \textbf{Transparenz.} Wenn ein Fehler unsererseits passiert ist, das ehrlich benennen.
  \item \textbf{Kein Streit.} Auch bei ungerechten Bewertungen freundlich antworten; eine ruhige Antwort wirkt souveräner als ein Gegenargument.
\end{itemize}

\subsection{Manipulation: warum sie sich nicht lohnt}

Gekaufte Bewertungen, fake Reviews oder das Drücken negativer Bewertungen funktionieren kurzfristig -- mittelfristig sanktionieren Plattformen rigoros (Sichtbarkeitsverlust, Account-Sperre). Außerdem verlieren Sie das Frühwarnsignal: wenn Sie nicht wissen, was wirklich schief läuft, können Sie es nicht beheben.

\begin{merkbox}[Bewertungen als Frühwarnsystem]
Bewertungen sind nicht nur ein \emph{Verkaufssignal}, sondern ein \emph{Prozesssignal}. Wenn drei Wochen lang Kritik an Lieferzeiten häuft, ist das ein Logistikproblem -- und kein Bewertungsproblem. Wer Bewertungen nicht als Frühwarnsystem nutzt, verliert ein zentrales Steuerungsinstrument.
\end{merkbox}

% --- 7. INTEGRIERTES SYSTEM ----------------------------------------------
\section{SEO, Content, Retention, Bewertung -- ein System}

Die vier Hebel dieses Tages -- SEO, Content, Retention, Bewertung -- sind keine voneinander getrennten Disziplinen. Sie verstärken sich gegenseitig:

\subsection{Vier Verbindungen}

\begin{enumerate}
  \item \textbf{SEO bringt Klicks -- Content macht sie zum Kauf.} Ohne saubere Produktseiten verpufft jede Sichtbarkeit.
  \item \textbf{Content erzeugt korrekte Erwartung -- Erwartung gibt gute Bewertungen.} Wer überverspricht, erntet schlechte Bewertungen.
  \item \textbf{Bewertungen treiben Ranking -- Ranking treibt Sichtbarkeit.} Ein Kreislauf, der bei jeder Stelle gebrochen werden kann.
  \item \textbf{Retention reduziert Akquisitionskosten -- niedrige CAC erlaubt höhere Marge.} Wer Bestandskunden hält, kann sich Wettbewerb auf neuen Suchen leisten.
\end{enumerate}

\subsection{Was eine Plattform-Performance-Architektur leistet}

\begin{itemize}
  \item Sie definiert pro Produktgruppe eine Keyword-Strategie (inkl.\ bewusst ausgeschlossener Keywords).
  \item Sie spezifiziert Content-Standards für Titel, Beschreibung, Bilder, FAQ.
  \item Sie hat einen Retention-Mechanismus für Zubehör, Bundles und Wiederkauf.
  \item Sie betreibt ein aktives Bewertungsmanagement mit Reaktionsregeln und Lernschleifen.
  \item Sie misst KPIs auf allen vier Achsen -- nicht nur Umsatz.
\end{itemize}

\quelle{vgl.\ Kotler et al., 2017; Homburg, 2020}

\begin{eselsbox}[Eselsbrücke: \akzent{S-C-R-B} -- vier Hebel im System]
Vier Buchstaben, vier Hebel:\\[2pt]
\textbf{S}EO -- werden wir gefunden?\\
\textbf{C}ontent -- macht die Produktseite den Klick zum Kauf?\\
\textbf{R}etention -- kommt der Kunde wieder?\\
\textbf{B}ewertungen -- vertrauen neue Kunden uns?\\[2pt]
\emph{Wer einen der vier vernachlässigt, gewinnt nie das Ganze.}
\end{eselsbox}

% --- 8. FALLSTUDIE-LEITFADEN ---------------------------------------------
\section{Fallstudie-Leitfaden: ToolMaster Pro}

In den Unterrichtseinheiten 4--5 bearbeiten Sie die Fallstudie ``ToolMaster Pro'' -- ein mittelständischer Anbieter hochwertiger Elektrowerkzeuge mit 2{,}8\,Mio.\,EUR Plattformumsatz, 24\,\% Marge und einem Bewertungsdurchschnitt von 4{,}3 Sternen. Wettbewerber sind günstiger und sichtbarer.

\subsection{Analytischer Aufbau}

\begin{enumerate}
  \item \textbf{Ursachen der schwachen Performance.} Wo liegt das Hauptproblem -- Content, Sichtbarkeit, Bewertungen, Retention?
  \item \textbf{Keyword-Strategie.} Welche Hauptkeywords, Long-Tail-, Problem-Keywords? Welche bewusst nicht (``billig'')?
  \item \textbf{Produkttitel und Beschreibung optimieren.} Beispiel-Optimierung an einer Produktseite.
  \item \textbf{Content-Maßnahmen.} Bilder, Vergleichstabellen, FAQ, Anwendungsvideos.
  \item \textbf{Retention.} Wie nutzen wir Zubehör- und Ersatzakku-Wiederkauf?
  \item \textbf{Bewertungsmanagement.} Welche vier Aufgaben (Ermöglichen, Analysieren, Beantworten, Lernen)?
  \item \textbf{Budgetverteilung.} Wie verteilen sich 90.000\,EUR auf 6 Monate?
  \item \textbf{Entscheidung unter Unsicherheit.} Warum \emph{nicht} Preisaktionen?
\end{enumerate}

\subsection{Erwartete Empfehlungsrichtung}

In einer typischen Musterlösung wird \emph{nicht} auf aggressive Preissenkungen gesetzt, sondern auf Content-Qualität, Keyword-Strategie und Bewertungsmanagement. Das Budget verteilt sich auf Keyword-/Content-Optimierung, neue Bilder/Videos, FAQ und Kompatibilitätsdaten, Bewertungsmanagement, Retention/Bundle-Maßnahmen und ein Testbudget. Die Entscheidung unter Unsicherheit: ``Trotz Wettbewerbsdruck verzichten wir auf pauschale Preisaktionen, weil sie die Qualitätspositionierung beschädigen und mittelfristig die Marge zerstören -- stattdessen testen wir Bundles und Content-Optimierung an ausgewählten Produktgruppen.''

\begin{merkbox}[Argumentationsmuster für die Fallstudie]
Eine gute Plattform-Performance-Strategie enthält:
\begin{enumerate}
  \item \textbf{Welche Keywords -- inklusive welche bewusst nicht.}
  \item \textbf{Welche Content-Maßnahmen mit welchem Kundennutzen.}
  \item \textbf{Welcher Retention-Mechanismus mit welcher KPI.}
  \item \textbf{Welches Bewertungsmanagement mit welcher Reaktionsregel.}
  \item \textbf{Welche Verzichtsentscheidung gegen kurzfristige Reichweite oder Preissenkung.}
\end{enumerate}
\end{merkbox}

% --- 9. MANAGEMENT-PERSPEKTIVE ------------------------------------------
\section{Management-Perspektive: Zielkonflikte}

Auf Wissens- und Anwendungsebene sieht eine gute Plattform-Strategie wie eine sauber priorisierte Liste aus. Auf Managementebene geht es um Zielkonflikte, die nicht ``richtig gelöst'' werden, sondern \emph{entschieden} werden müssen. \quelle{vgl.\ Homburg, 2020; Kotler et al., 2017}

\subsection{Fünf zentrale Zielkonflikte}

\begin{enumerate}
  \item \textbf{Reichweite vs.\ Markenwirkung.} Maximale Sichtbarkeit kann auf Kosten der Positionierung gehen -- besonders, wenn man Generika-Keywords gegen Discounter optimiert.
  \item \textbf{Conversion vs.\ Marge.} Aggressive Preisaktionen erhöhen Conversion -- senken aber Marge und beschädigen Preisanker.
  \item \textbf{Kurzfristiger Umsatz vs.\ langfristige Plattform-Performance.} Eine Werbekampagne bringt kurzfristig Klicks; eine saubere Content-Strategie wirkt erst nach Wochen, aber dauerhaft.
  \item \textbf{Kundenerwartung vs.\ Markenversprechen.} Schnelle Lieferzeiten und 24h-Versprechen sind verlockend -- aber wer sie nicht hält, verliert in Bewertungen.
  \item \textbf{Eigene Sichtbarkeit vs.\ Plattformabhängigkeit.} Wer ausschließlich auf einer Plattform wächst, ist 100\,\% abhängig. Diversifikation kostet -- aber sie reduziert das strategische Risiko.
\end{enumerate}

\subsection{Plattform-Performance ist Senior-Level-Thema}

Plattform-Verkauf wirkt operativ -- Titel, Bilder, FAQ, Bewertungen. Diese Operativität täuscht darüber hinweg, dass die Entscheidungen \emph{Senior-Level-Charakter} haben: Welche Marke kommunizieren wir? Welche Kundenerwartung erzeugen wir? Welchen Preisanker setzen wir? Welche Plattform­abhängigkeit lassen wir entstehen? Wer diese Fragen an die Plattform-Operativ-Abteilung delegiert, delegiert die Markenführung -- und das kann Senior-Management nicht.

\subsection{Senior-Level-Test}

Eine senior-level-fähige Plattform-Strategie lässt sich in einem Satz erklären, der eine Verzichtsentscheidung enthält. Beispiel: ``Wir optimieren auf Qualitäts-Keywords und Bundles -- wir verzichten bewusst auf 'billig'-Keywords und pauschale Rabattaktionen, weil sie unsere Marke beschädigen und mittelfristig die Marge zerstören.''

\begin{eselsbox}[Eselsbrücke: \akzent{R-C-K-E-A}]
Fünf Zielkonflikte, die im Plattform-Verkauf auf Senior-Level entschieden werden:\\[2pt]
\textbf{R}eichweite vs.\ Markenwirkung.\\
\textbf{C}onversion vs.\ Marge.\\
\textbf{K}urzfrist vs.\ Plattform-Performance.\\
\textbf{E}rwartung vs.\ Markenversprechen.\\
\textbf{A}bhängigkeit (Plattform) vs.\ Diversifikation.\\[2pt]
\emph{Fünf Konflikte, die kein Plattform-Operativ-Team allein verantworten sollte.}
\end{eselsbox}

\begin{merkbox}[Schluss-Merksatz für Tag 7]
Plattform-Verkauf wirkt operativ, ist aber eine \textbf{strategische Disziplin}. Sichtbarkeit, Content, Retention und Bewertungen sind kein Werbethema, sondern ein zusammenhängendes Steuerungssystem für Margen, Marke und Verteidigungsfähigkeit. \emph{Wer Plattform-Performance nicht als Architektur denkt, baut Klickkampagnen ohne Plan.}
\end{merkbox}

% --- 10. LITERATUR -------------------------------------------------------
\clearpage
\section{Literatur}

Im Skript verwendete und für die Vertiefung empfohlene Quellen. Zitierstil: APA-7.

\begin{description}[style=nextline,leftmargin=18pt,labelindent=0pt,itemsep=4pt]
  \item[Chevalier, J.\,A., \& Mayzlin, D. (2006).] The effect of word of mouth on sales: Online book reviews. \emph{Journal of Marketing Research, 43}(3), 345--354. \href{https://doi.org/10.1509/jmkr.43.3.345}{https://doi.org/10.1509/jmkr.43.3.345}
  \item[Ellis, S., \& Brown, M. (2017).] \emph{Hacking growth: How today's fastest-growing companies drive breakout success}. Crown Business.
  \item[Enge, E., Spencer, S., \& Stricchiola, J.\,C. (2015).] \emph{The art of SEO: Mastering search engine optimization} (3.\ Aufl.). O'Reilly.
  \item[Fishkin, R. (2018).] \emph{Lost and founder: A painfully honest field guide to the startup world}. Portfolio.
  \item[Homburg, C. (2020).] \emph{Marketingmanagement: Strategie -- Instrumente -- Umsetzung -- Unternehmensführung} (7.\ Aufl.). Springer Gabler.
  \item[Kotler, P., Kartajaya, H., \& Setiawan, I. (2017).] \emph{Marketing 4.0: Moving from traditional to digital}. Wiley.
  \item[Lemon, K.\,N., \& Verhoef, P.\,C. (2016).] Understanding customer experience throughout the customer journey. \emph{Journal of Marketing, 80}(6), 69--96. \href{https://doi.org/10.1509/jm.15.0420}{https://doi.org/10.1509/jm.15.0420}
  \item[Reichheld, F.\,F. (2003).] The one number you need to grow. \emph{Harvard Business Review, 81}(12), 46--54.
\end{description}

% --- 11. GLOSSAR ---------------------------------------------------------
\clearpage
\section{Glossar}

Die wichtigsten Begriffe des Tages -- kurz definiert, in alphabetischer Reihenfolge.

\begin{description}[style=nextline,leftmargin=0pt,labelindent=0pt,itemsep=4pt]
  \item[Bewertung] Quantitative oder qualitative Rückmeldung eines Kunden zu Produkt, Service, Lieferung; Vertrauenssignal und Ranking-Faktor.
  \item[Bewertungsmanagement] Strukturierter Prozess aus Ermöglichen, Analysieren, Beantworten und Lernen aus Bewertungen.
  \item[Bundle] Kombination aus Hauptprodukt, Zubehör und/oder Service zu einem Festpreis; reduziert Vergleichbarkeit, erhöht Warenkorbwert.
  \item[Content-Optimierung] Verbesserung von Produkttitel, Beschreibung, Bildern, FAQ und Vergleichsinhalten zur Steigerung von Vertrauen und Conversion.
  \item[Conversion Rate] Anteil Klicks auf eine Produktseite, die zum Kauf führen.
  \item[Customer Acquisition Cost (CAC)] Gesamtkosten zur Gewinnung eines neuen Kunden.
  \item[Frühwarnsystem (Bewertung)] Nutzung wiederkehrender Bewertungsmuster als Indikator für Prozess- oder Produktprobleme.
  \item[Hauptkeyword] Generisches Suchwort mit hohem Volumen und hohem Wettbewerb.
  \item[Kaufnahes Keyword] Suchbegriff mit erkennbarer Kaufabsicht.
  \item[Keyword] Suchbegriff oder -phrase, mit der Kunden ein Bedürfnis ausdrücken.
  \item[Keyword-Strategie] Geordnete Auswahl von Suchbegriffen pro Produktgruppe, einschließlich bewusst ausgeschlossener Begriffe.
  \item[Klickrate (CTR)] Anteil Sichtungen, die zu einem Klick führen.
  \item[Long-Tail-Keyword] Längere, präzisere Suchphrase mit niedrigerem Volumen, aber höherer Kaufabsicht.
  \item[Markenpassung (Keyword)] Eigenschaft eines Keywords, die zur Positionierung des Anbieters passt -- oder eben nicht.
  \item[Plattform-SEO] Suchmaschinenoptimierung innerhalb einer Plattform (z.\,B.\ Marktplatz), mit eigenen Ranking-Faktoren und Algorithmen.
  \item[Problemorientiertes Keyword] Suchbegriff, der eine Anwendungssituation beschreibt -- z.\,B.\ ``Akkuschrauber für Dauereinsatz''.
  \item[Produkttitel] Wichtigste einzelne Optimierungsstelle auf Plattformen; verbindet Produkttyp, Hauptkeyword, Differenzierung und Anwendungsbezug.
  \item[Retention] Kundenbindung über Wiederkauf, Erweiterung, Abonnement, Empfehlung.
  \item[Retourenquote] Anteil verkaufter Produkte, die zurückgeschickt werden; senkt Plattform-Sichtbarkeit.
  \item[Sichtbarkeit] Wahrscheinlichkeit, dass ein Produkt in relevanten Suchen angezeigt wird; Ergebnis aus Algorithmus, Performance und Bewertungen.
  \item[Suchintention] Absicht, die ein Nutzer mit seiner Suche verfolgt: informieren, vergleichen, kaufen, wiederbestellen.
  \item[Verkäuferperformance] Bewertungsmaßstab der Plattform für Anbieter -- Reklamationsquote, Antwortzeiten, Konfliktlösung, Lieferqualität.
  \item[Vergleichstabelle] Inhalt, der Unterschiede zwischen eigenen Produktvarianten transparent macht.
  \item[Wiederkaufrate] Anteil Kunden, die innerhalb eines definierten Zeitraums erneut beim Anbieter kaufen.
  \item[Zielkonflikt (Plattform-Verkauf)] Spannungsfeld zwischen Reichweite, Marge, Marke, Kundenerwartung und Plattform­abhängigkeit -- bewusst entscheidbar, nicht ``lösbar''.
\end{description}

\vspace{1cm}
\noindent{\color{lpAccent}\rule{\linewidth}{0.6pt}}\\[4pt]
{\footnotesize\sffamily\color{lpGray}Lernplatt \textperiodcentered{} by Can Siebert \textperiodcentered{} Kurs 22 (Bo 143) \textperiodcentered{} Tag 7 \textperiodcentered{} \textcopyright{} 2026}

\end{document}
